\documentclass[11pt]{article}
\usepackage[margin=1in]{geometry}
\usepackage{amsmath,amssymb,booktabs,authblk,url}
\newcommand{\grad}{\nabla}
\newcommand{\half}{\tfrac12}
\newcommand{\eps}{\varepsilon}

\title{The Clay Navier--Stokes Problem:\\[0.3em]
\large Formulation and Solution}
\author[1]{Claes Johnson}
\affil[1]{KTH Royal Institute of Technology, SE-100 44 Stockholm, Sweden \\
\texttt{claesjohnson@gmail.com}}
\date{}

\begin{document}
\maketitle

\begin{abstract}
\noindent The official formulation of the Clay Millennium Problem on the Navier--Stokes equations
asks whether smooth initial data give a solution that stays smooth for all time, or whether the
velocity can blow up in finite time. The alternative is exhaustive only if those are the two
possibilities. They are not. The solutions of physical interest are \emph{turbulent}: not smooth,
because velocity gradients grow as $\nu^{-1/2}$ and are enormous at the Reynolds numbers of any
real flow, and not singular, because the velocity itself stays bounded. A dichotomy that omits the
only case that matters cannot be answered in a way that says anything about fluids, and the case it
does cover --- persistent smoothness --- is either trivial or false depending on how it is read.

This article argues that the problem as posed is mis-posed, and states what should replace it.
The replacement is Hadamard's criterion: not global smoothness, but \emph{wellposedness}. A
solution concept earns physical meaning when the quantities computed from it are stable under
perturbation. For turbulent solutions those quantities are not pointwise velocities, which are
indeterminate, but mean-value outputs such as drag and lift, which are not. We set out the
distinction, the evidence that turbulent solutions are wellposed in outputs while smooth solutions
are not wellposed at all, and the sense in which computation --- finite element solution with
residual-based stabilisation, error control by duality --- constitutes an answer. The reformulated
problem has a solution, and the solution is that turbulence is what Navier--Stokes has to say.
\end{abstract}

\section{The official formulation}

The Clay Mathematics Institute problem, as set out by Fefferman \cite{fefferman}, concerns the
incompressible Navier--Stokes equations
\begin{equation}\label{eq:ns}
\dot u + (u\cdot\grad)u - \nu\Delta u + \grad p = f,\qquad \grad\cdot u = 0,
\end{equation}
on $\mathbb{R}^3$ or a periodic box, with $u(\cdot,0)=u^0$ smooth and divergence free, and with the
viscosity $\nu$ normalised to one. Four statements are offered, of which a proof of any one settles
the prize: existence and smoothness on $\mathbb{R}^3$; the same with periodic data; or a
counterexample to either --- a solution which is initially smooth and whose velocity becomes
unbounded in finite time.

The structure of the question is therefore a dichotomy. Either the solution is smooth for all time,
or it is not, and the only stated way for it not to be is blowup. \emph{Tertium non datur.}

\section{Why the dichotomy is empty}

\paragraph{The physical solution is neither.} A fluid at a Reynolds number of interest --- $10^6$
and upward for an aircraft, a ship, a weather system, a pipe --- is turbulent. Turbulent flow is not
smooth in any useful sense: the velocity varies on a length scale that shrinks with viscosity, and
the gradients scale as
\begin{equation}\label{eq:grad}
|\grad u| \sim \nu^{-1/2},
\end{equation}
so that at $\nu=10^{-6}$ they are of order $10^{3}$ times the velocity scale, and at the viscosity of
air over a wing, larger still. Nor is it singular: the velocity stays bounded, the kinetic energy
stays finite, and nothing goes to infinity anywhere. Kolmogorov's scaling, in the form Onsager gave
it, puts the velocity field at H\"older continuity $1/3$ --- continuous, with unbounded derivatives
in the limit. That is precisely the case the official alternative does not contain.

\paragraph{The smooth case is trivial or false.} If smoothness is to be read as ``the solution
exists and is smooth for a short time'', the statement is classical and was not worth a prize. If it
is to be read as ``the solution stays smooth for all time'', then for the flows of interest it is
false, and not because of blowup: a smooth solution of \eqref{eq:ns} at small viscosity is
\emph{unstable}. Perturb potential flow around a bluff body and it does not stay potential flow; it
develops rotational separation and becomes turbulent. The smooth solution exists as a formal object
and is not observed, because nothing keeps a flow in it. So the first horn of the dichotomy
describes solutions that are mathematically available and physically absent, and the second
describes solutions that do not occur.

\paragraph{The normalisation removes the phenomenon.} Setting $\nu=1$ is harmless if one is asking
about existence in the abstract and fatal if one is asking about fluids. All of the structure in
\eqref{eq:grad} is structure in $\nu$. A formulation that normalises the viscosity away has
normalised away turbulence, which is to say it has normalised away the subject.

\paragraph{The quantitative results confirm rather than relieve this.} Work in the analytic
tradition has produced conditional regularity bounds of the form: if the velocity is bounded by $A$,
then the derivatives are bounded by an expression exponential in $A$ iterated several times. As a
theorem this is real. As a statement about a fluid it is empty: a bound of the size
$\exp\exp\exp A$ is not a bound that any flow can be said to respect, and no measurement could
distinguish it from no bound at all. Tartar's objection to the problem's framing --- that the
selected questions have almost no physical content, and appear to have been chosen by specialists in
harmonic analysis rather than by anyone studying continuum mechanics --- is, on this evidence, well
taken \cite{tartar}.

\section{What should replace it: wellposedness}

The defect in the official formulation is not that it asks a hard question. It is that it asks about
the wrong property. Smoothness is not what makes a solution concept physically meaningful.
Wellposedness is.

Hadamard's criterion has three parts: a solution exists, it is unique, and it depends continuously
on the data. The third is the one that carries the physical content, because data are never exact.
A problem whose solution is not stable under perturbation of the data is not a model of anything
measurable, whatever its smoothness.

\paragraph{Weak solutions exist.} Leray showed in 1934 that weak solutions of \eqref{eq:ns} exist
globally in time \cite{leray}. What was left open was not existence but whether the weak solution
concept is \emph{wellposed} --- whether it determines anything. That is the question worth asking,
and it is not the question the Clay problem asks.

\paragraph{Wellposedness is a property of outputs, not of solutions.} Here is the point on which the
reformulation turns. For a turbulent solution, the pointwise velocity at a given place and time is
not determined: perturb the data in the last decimal and the pointwise field is different
everywhere within a short time. In that sense the solution is not wellposed and never will be.

But nobody measures the pointwise velocity of a turbulent flow, and no engineering question depends
on it. What is measured, and what is asked for, are \emph{mean-value outputs}: the drag on a body,
the lift on a wing, a pressure averaged over a surface, a flux through a section. These are
functionals of the solution, and the question of wellposedness must be asked of them:
\begin{equation}\label{eq:output}
M(u) = \int_Q u\cdot\psi\,dx\,dt
\end{equation}
for a given weight $\psi$. The claim is that $M(u)$ is stable for turbulent solutions even though
$u$ is not.

\paragraph{The mechanism, and why it inverts the classical expectation.} Stability of an output is
governed by an associated dual, or adjoint, problem: linearise about the computed solution, run it
backward from the output functional, and the growth of the dual solution measures how strongly a
perturbation of the data or of the residual can move the output. What the dual problem shows is the
reverse of what intuition suggests. Over a \emph{turbulent} velocity field --- oscillating rapidly in
space and time --- the dual solution grows only modestly, because the linearised coefficients
oscillate and their effect largely cancels. Over a \emph{smooth} field the dual solution grows,
because there is nothing to cancel. Turbulence stabilises the outputs that smoothness destabilises.

This is the sense in which the reformulated problem has an answer, and the answer is not the one the
dichotomy anticipates. Smooth solutions are not wellposed. Turbulent solutions are wellposed in
mean-value outputs. The physically meaningful solution concept is the one the official formulation
treats as the failure case.

\section{The reformulated problem}

We state the replacement plainly.

\begin{quote}
\textbf{Problem.} For the incompressible Navier--Stokes equations at small viscosity, with smooth
divergence-free initial data, show that the weak solution is \emph{wellposed in mean-value outputs}:
that functionals of the form \eqref{eq:output} depend continuously on the data, with quantitative
error control, even though the pointwise solution does not. Equivalently: establish the second law
of thermodynamics for fluids --- that a turbulent solution dissipates kinetic energy at a rate which
does not vanish with the viscosity, and that this dissipation is what makes the outputs stable.
\end{quote}

Two features of this statement are worth marking. It contains the viscosity, so it is a statement
about fluids and not about a normalised abstraction. And it asks for stability rather than
regularity, so it can be answered affirmatively for the solutions that actually occur instead of
being answered vacuously for solutions that do not.

\section{In what sense this constitutes a solution}

\paragraph{The solutions are constructed, not assumed.} A finite element method for \eqref{eq:ns}
with residual-based stabilisation --- Galerkin least squares on a mesh, piecewise polynomial
velocity and pressure, no turbulence model supplied by the user --- produces solutions which become
turbulent even when initiated from smooth potential data, and which stay bounded. The stabilisation
supplies dissipation in proportion to the residual, so the dissipation is generated by the
computation rather than prescribed by a model. In the vanishing-viscosity limit the same scheme is
an Euler solver whose turbulent dissipation appears automatically. The non-smooth,
non-singular third case is therefore not a conjecture: it is what the equations produce when they
are solved.

\paragraph{The residual behaves as the theory requires.} For such solutions the Navier--Stokes
residual $R(U,P)$ is small in a weak norm and large in a strong one. Small weakly is what makes the
computed field a solution; large strongly is the statement that it is not smooth. A method whose
residual were small in both would be computing something that does not exist.

\paragraph{The error is controlled by duality, not by asserted convergence.} An a posteriori error
estimate of the form
\begin{equation}\label{eq:apost}
|M(u) - M(U)| \;\le\; \int_Q R(U,P)\cdot\varphi\,dx\,dt \;\le\; S\,\|R(U,P)\|_{-1},
\end{equation}
with $\varphi$ the dual solution and $S$ its stability factor, converts a computed solution into a
statement with a bound attached. The stability factor is not assumed; it is computed by solving the
dual problem. It is through \eqref{eq:apost} that the wellposedness claim becomes quantitative: $S$
is moderate for mean-value outputs over turbulent solutions and large for pointwise outputs and for
smooth solutions, which is the distinction of \S3 expressed as a number.

\paragraph{Computation as proof.} That an answer rests on computation is not a defect peculiar to
this problem. The four-colour theorem was settled the same way, and the settlement was accepted
because the computation was reproducible and its error controlled. The same standard applies here,
and \eqref{eq:apost} is what meets it: the claim is not that a particular simulation looked
turbulent, but that the output of the computation comes with a bound which can be evaluated, and
which is small.

\section{What the answer says about the original question}

To the official question --- do smooth data give solutions that remain smooth for all time --- the
answer implied by the above is no, and not by blowup. Smooth data give solutions that become
turbulent: non-smooth, with gradients growing as $\nu^{-1/2}$, bounded in velocity, dissipating at a
rate which does not vanish as $\nu\to0$. Blowup, the only failure mode the official formulation
admits, does not occur and is not the interesting question. The interesting question is the one the
formulation leaves out.

This does not claim the prize, and it is not intended to. It claims something less procedural and
more useful: that the problem as posed cannot be answered in a way that illuminates fluid mechanics,
because the alternative it offers is not exhaustive over the cases that occur; and that the
adjacent, answerable question --- wellposedness of mean-value outputs of turbulent weak solutions
--- is both the physically meaningful one and one to which the answer is affirmative and
computable.

\begin{thebibliography}{9}
\bibitem{fefferman} C.~Fefferman, \emph{Existence and smoothness of the Navier--Stokes equation},
Official Problem Description, Clay Mathematics Institute, 2000.
\bibitem{leray} J.~Leray, \emph{Sur le mouvement d'un liquide visqueux emplissant l'espace},
Acta Math.\ \textbf{63} (1934) 193--248.
\bibitem{tartar} L.~Tartar, remarks on the selection of the Millennium Problems, quoted in
discussion of the Clay problem formulation.
\bibitem{cfd} J.~Hoffman and C.~Johnson, \emph{Computational Turbulent Incompressible Flow},
Applied Mathematics: Body and Soul, vol.~4, Springer, 2007.
\bibitem{blog} C.~Johnson, posts under the label \emph{clay problem},
\url{https://claesjohnson.blogspot.com/search/label/clay\%20problem}.
\end{thebibliography}

\end{document}
